\chapter{2020年新高考II卷}

\section{单选题}

\begin{problem}
设集合 \(A=\{2,3,5,7\}\),\(B=\{1,2,3,5,8\}\),则 \(A\cap B=\)(\quad)

\choices
    {\(\{1,8\}\)}
    {\(\{2,5\}\)}
    {\(\{2,3,5\}\)}
    {\(\{1,2,3,5,8\}\)}
\end{problem}

\begin{problem}
\((1+2\i)(2+\i)=\)(\quad)

\choices
    {\(-5\i\)}
    {\(5\i\)}
    {\(-5\)}
    {\(5\)}
\end{problem}

\begin{problem}
在 \(\triangle ABC\) 中,\(D\) 是 \(AB\) 边上的中点,则 \(\overrightarrow{CB}=\)(\quad)

\choices
    {\(2\overrightarrow{CD}-\overrightarrow{CA}\)}
    {\(2\overrightarrow{CA}-\overrightarrow{CD}\)}
    {\(2\overrightarrow{CD}+\overrightarrow{CA}\)}
    {\(2\overrightarrow{CA}+\overrightarrow{CD}\)}
\end{problem}

\begin{problem}
日晷是中国古代用来测定时间的仪器,利用与晷面垂直的晷针投射到晷面的影子来测定时间. 把地球看成一个球(球心记为 \(O\)),地球上一点 \(A\) 的纬度是指 \(OA\) 与地球赤道所在平面所成角,点 \(A\) 处的水平面是指过点 \(A\) 且与 \(OA\) 垂直的平面. 在点 \(A\) 处放置一个日晷,若晷面与赤道所在平面平行,点 \(A\) 处的纬度为北纬 \(40^\circ\),则晷针与点 \(A\) 处的水平面所成角为(\quad)

\begin{center}
\bitmapfigure[float=inline,max width=0.42\linewidth]{img/2020/new_gaokao_paper_2/q4_fig1.png}
\end{center}

\choices
    {\(20^\circ\)}
    {\(40^\circ\)}
    {\(50^\circ\)}
    {\(90^\circ\)}
\end{problem}

\begin{problem}
某中学的学生积极参加体育锻炼,其中有 \(96\%\) 的学生喜欢足球或游泳,\(60\%\) 的学生喜欢足球,\(82\%\) 的学生喜欢游泳,则该中学既喜欢足球又喜欢游泳的学生数占该校学生总数的比例是(\quad)

\choices
    {62\%}
    {56\%}
    {46\%}
    {42\%}
\end{problem}

\begin{problem}
\(3\) 名大学生利用假期到 \(2\) 个山村参加扶贫工作,每名大学生只去 \(1\) 个村,每个村至少 \(1\) 人,则不同的分配方案共有(\quad)

\choices
    {\(4\) 种}
    {\(5\) 种}
    {\(6\) 种}
    {\(8\) 种}
\end{problem}

\begin{problem}
已知函数 \(f(x)=\lg(x^2-4x-5)\) 在 \((a,+\infty)\) 单调递增,则 \(a\) 的取值范围是(\quad)

\choices
    {\((-\infty,-1]\)}
    {\((-\infty,2]\)}
    {\([2,+\infty)\)}
    {\([5,+\infty)\)}
\end{problem}

\begin{problem}
若定义在 \(\R\) 的奇函数 \(f(x)\) 在 \((-\infty,0)\) 单调递减,且 \(f(2)=0\),则满足 \(xf(x-1)\geqslant0\) 的 \(x\) 的取值范围是(\quad)

\choices
    {\([-1,1]\cup[3,+\infty)\)}
    {\([-3,-1]\cup[0,1]\)}
    {\([-1,0]\cup[1,+\infty)\)}
    {\([-1,0]\cup[1,3]\)}
\end{problem}

\section{多选题}

\begin{problem}
我国新冠肺炎疫情进入常态化,各地有序推进复工复产,下面是某地连续 \(11\) 天复工复产指数折线图,下列说法正确的是(\quad)

\begin{center}
\bitmapfigure[float=inline,max width=0.88\linewidth]{img/2020/new_gaokao_paper_2/q9_fig1.png}
\end{center}

\choices
    {这 \(11\) 天复工指数和复产指数均逐日增加}
    {这 \(11\) 天期间,复产指数增量大于复工指数的增量}
    {第 \(3\) 天至第 \(11\) 天复工复产指数均超过 \(80\%\)}
    {第 \(9\) 天至第 \(11\) 天复产指数增量大于复工指数的增量}
\end{problem}

\begin{problem}
已知曲线 \(C:mx^2+ny^2=1\),则下列正确的是(\quad)

\choices
    {若 \(m>n>0\),则 \(C\) 是椭圆,其焦点在 \(y\) 轴上}
    {若 \(m=n>0\),则 \(C\) 是圆,其半径为 \(\sqrt n\)}
    {若 \(mn<0\),则 \(C\) 是双曲线,其渐近线方程为 \(y=\pm\sqrt{-\dfrac{m}{n}}x\)}
    {若 \(m=0\),\(n>0\),则 \(C\) 是两条直线}
\end{problem}

\begin{problem}
下图是函数 \(y=\sin(\omega x+\varphi)\) 的部分图像,则 \(\sin(\omega x+\varphi)=\)(\quad)

\begin{center}
\bitmapfigure[float=inline,max width=0.42\linewidth]{img/2020/new_gaokao_paper_2/q11_fig1.png}
\end{center}

\choices
    {\(\sin\left(x+\dfrac{\pi}{3}\right)\)}
    {\(\sin\left(\dfrac{\pi}{3}-2x\right)\)}
    {\(\cos\left(2x+\dfrac{\pi}{6}\right)\)}
    {\(\cos\left(\dfrac{5\pi}{6}-2x\right)\)}
\end{problem}

\begin{problem}
已知 \(a>0\),\(b>0\),且 \(a+b=1\),则(\quad)

\choices
    {\(a^2+b^2\geqslant\dfrac12\)}
    {\(2^{a-b}>\dfrac12\)}
    {\(\log_2a+\log_2b\geqslant-2\)}
    {\(\sqrt a+\sqrt b\leqslant\sqrt2\)}
\end{problem}

\section{填空题}

\begin{problem}
棱长为 \(2\) 的正方体 \(ABCD-A_1B_1C_1D_1\) 中,\(M\),\(N\) 分别为棱 \(BB_1\),\(AB\) 的中点,则三棱锥 \(A-NMD_1\) 的体积为\fillinblank{}.
\end{problem}

\begin{problem}
斜率为 \(\sqrt3\) 的直线过抛物线 \(C:y^2=4x\) 的焦点,且与 \(C\) 交于 \(A\),\(B\) 两点,则 \(|AB|=\fillinblank{}\).
\end{problem}

\begin{problem}
将数列 \(\{2n-1\}\) 与 \(\{3n-2\}\) 的公共项从小到大排列得到数列 \(\{a_n\}\),则 \(\{a_n\}\) 的前 \(n\) 项和为\fillinblank{}.
\end{problem}

\begin{problem}
某中学开展劳动实习,学生加工制作零件,零件的截面如图所示. \(O\) 为圆孔及轮廓圆弧 \(AB\) 所在圆的圆心,\(A\) 是圆弧 \(AB\) 与直线 \(AG\) 的切点,\(B\) 是圆弧 \(AB\) 与直线 \(BC\) 的切点,四边形 \(DEFG\) 为矩形,\(BC\perp DG\),垂足为 \(C\),\(\tan\angle ODC=\dfrac35\),\(BH\parallel DG\),\(EF=12\text{ cm}\),\(DE=2\text{ cm}\),\(A\) 到直线 \(DE\) 和 \(EF\) 的距离均为 \(7\text{ cm}\),圆孔半径为 \(1\text{ cm}\),则图中阴影部分的面积为\fillinblank{}\ \(\text{cm}^2\).

\begin{center}
\bitmapfigure[float=inline,max width=0.88\linewidth]{img/2020/new_gaokao_paper_2/q16_fig1.png}
\end{center}
\end{problem}

\section{解答题}

\begin{problem}
在
\begin{circlelist}
\item \(ac=\sqrt3\),
\item \(c\sin A=3\),
\item \(c=\sqrt3 b\)
\end{circlelist}
这三个条件中任选一个,补充在下面问题中,若问题中的三角形存在,求 \(c\) 的值;若问题中的三角形不存在,说明理由.

问题:是否存在 \(\triangle ABC\),它的内角 \(A\),\(B\),\(C\) 的对边分别为 \(a\),\(b\),\(c\),且 \(\sin A=\sqrt3\sin B\),\(C=\dfrac{\pi}{6}\),\fillinblank{}?
\end{problem}

\begin{problem}
已知公比大于 \(1\) 的等比数列 \(\{a_n\}\) 满足 \(a_2+a_4=20\),\(a_3=8\).

\begin{enumerate}
\item 求 \(\{a_n\}\) 的通项公式;
\item 求 \(a_1a_2-a_2a_3+\cdots+(-1)^{n-1}a_na_{n+1}\).
\end{enumerate}
\end{problem}

\begin{problem}
为加强环境保护,治理空气污染,环境监测部门对某市空气质量进行调研,随机抽查了 \(100\) 天空气中的 \(\mathrm{PM}_{2.5}\) 和 \(\mathrm{SO}_2\) 浓度(单位:\(\mu\text{g}/\text{m}^3\)),得下表:

\begin{center}
\begin{tabular}{c|ccc}
\(\mathrm{PM}_{2.5}\backslash\mathrm{SO}_2\) & \([0,50]\) & \((50,150]\) & \((150,475]\) \\
\hline
\([0,35]\) & 32 & 18 & 4 \\
\((35,75]\) & 6 & 8 & 12 \\
\((75,115]\) & 3 & 7 & 10 \\
\end{tabular}
\end{center}

\begin{enumerate}
\item 估计事件``该市一天空气中 \(\mathrm{PM}_{2.5}\) 浓度不超过 \(75\),且 \(\mathrm{SO}_2\) 浓度不超过 \(150\)''的概率;
\item 根据所给数据,完成下面的 \(2\times2\) 列联表:

\begin{center}
\begin{tabular}{c|cc}
\(\mathrm{PM}_{2.5}\backslash\mathrm{SO}_2\) & \([0,150]\) & \((150,475]\) \\
\hline
\([0,75]\) & & \\
\((75,115]\) & & \\
\end{tabular}
\end{center}

\item 根据(2)中的列联表,判断是否有 \(99\%\) 的把握认为该市一天空气中 \(\mathrm{PM}_{2.5}\) 浓度与 \(\mathrm{SO}_2\) 浓度有关?
\end{enumerate}

附:\(K^2=\dfrac{n(ad-bc)^2}{(a+b)(c+d)(a+c)(b+d)}\).

\begin{center}
\begin{tabular}{c|ccc}
\(P(K^2\geqslant k)\) & \(0.050\) & \(0.010\) & \(0.001\) \\
\hline
\(k\) & \(3.841\) & \(6.635\) & \(10.828\)
\end{tabular}
\end{center}
\end{problem}

\begin{problem}
如图,四棱锥 \(P-ABCD\) 底面为正方形,\(PD\perp\) 底面 \(ABCD\). 设平面 \(PAD\) 与平面 \(PBC\) 的交线为 \(l\).

\begin{enumerate}
\item 证明:\(l\perp\) 平面 \(PDC\);
\item 已知 \(PD=AD=1\),\(Q\) 为 \(l\) 上的点,\(QB=\sqrt2\),求 \(PB\) 与平面 \(QCD\) 所成角的正弦值.
\end{enumerate}

\bitmapfigure[float=inline,max width=0.42\linewidth]{img/2020/new_gaokao_paper_2/q20_fig1.png}
\end{problem}

\begin{problem}
已知椭圆 \(C:\dfrac{x^2}{a^2}+\dfrac{y^2}{b^2}=1\)(\(a>b>0\))过点 \(M(2,3)\),点 \(A\) 为其左顶点,且 \(AM\) 的斜率为 \(\dfrac12\).

\begin{enumerate}
\item 求 \(C\) 的方程;
\item 点 \(N\) 为椭圆上任意一点,求 \(\triangle AMN\) 的面积的最大值.
\end{enumerate}
\end{problem}

\begin{problem}
已知函数 \(f(x)=ae^{x-1}-\ln x+\ln a\).

\begin{enumerate}
\item 当 \(a=e\) 时,求曲线 \(y=f(x)\) 在点 \((1,f(1))\) 处的切线与两坐标轴围成的三角形的面积;
\item 若 \(f(x)\geqslant1\),求 \(a\) 的取值范围.
\end{enumerate}
\end{problem}

